\section{Introduction}
Presentations communicate information to a particular audience in a particular setting. Requests for presentation generation not only specify a topic but often also involve requirements on slide count, formatting, mandatory content, and source grounding; generated slides must also avoid severe visual defects such as cropping and occlusion. A central question in evaluating these systems is whether their outputs satisfy these requirements. We refer to hard instruction failures, missing required content, and page-level cropping or occlusion collectively as \emph{blocking failures}, and focus on their joint diagnosis.

A high overall score does not imply the absence of failures that block practical use. We collect 240 outputs from six publicly available commercial systems. The six representative evaluation frameworks considered in Section~\ref{sec:screening} assign relatively high mean scores to the evaluated outputs: SlideAudit averages above .97, LecEval above .90, and PPTBench-style above .81. SlideGate identifies 124 decks with at least one blocking failure; Skywork has the most, at 27 of 40 tasks. In the initial grouping analysis, five of the six scoring schemes show no significant score difference between decks with and without detected blocking failures (Mann--Whitney $U$; Section~\ref{sec:screening}).

Consider an academic-talk task with a slide-count constraint. One system exceeds the allowed count, yet receives .993 from SlideAudit, .851 from PPTBench-style scoring, and .873 from LecEval. The violation thus coexists with relatively high scores under all three schemes. Content failures can arise similarly: a training task explicitly requires the second lesson plan as its source, but a system generates the deck entirely from the first lesson. A content evaluator without source materials cannot directly verify whether the specified source was used.

We introduce SlideGate, a presentation-generation benchmark comprising 40 tasks and outputs from six systems---AnyGen, Kimi, Doubao, Skywork, MiniMax, and Manus. Tasks span document attachments, editing instructions, template references, and other input forms. Through multiple independent generations and critic-model review, SlideGate converts each request and its parseable sources into a task-specific verification specification. A multimodal evaluator then checks each item, enabling systematic, dimension-specific assessment with traceable evidence.

Our contributions are:
\begin{enumerate}
\item \textbf{Evaluation protocol.} We organize evaluation around intended use and separately report instruction following, content quality, and visual presentation. Specifications are automatically derived from requests and parseable sources, frozen, and reused across all systems solving the same task. New tasks with supported input formats reuse the procedure with their own specifications.
\item \textbf{Benchmark.} SlideGate contains 40 tasks across seven scenarios and five input forms, with 2,191 verification items. The main experiment evaluates 240 outputs from six commercial presentation-generation systems.
\item \textbf{Findings and validation on existing systems.} Evaluating six commercial systems shows that blocking failures can coexist with high overall scores. SlideGate localizes these failures to specific items and pages, and validates the corresponding evaluation dimensions against programmatic references, independent human annotations, and designer preferences.
\end{enumerate}
